\hypertarget{chapter-05-page-117}{}
\subsection{向量值推广}
\label{sec:ch5-vector-valued-extension}

设 $B$ 是可分 Banach 空间. 若对每个 $b'\in B^*$, 标量函数 $x\mapsto\langle F(x),b'\rangle$ 都可测, 则称 $F:\mathbb R^n\to B$ 强可测. 此时 $x\mapsto\|F(x)\|_B$ 也可测. 定义 $L^p(B)$ 为满足
\[
 \|F\|_{L^p(B)}=\left(\int_{\mathbb R^n}\|F(x)\|_B^p\,dx\right)^{1/p}<\infty
\]
的可测函数等价类, 并按通常方式定义 $L^\infty(B)$. $L^p(B)$ 是 Banach 空间. 有限和 $F=\sum_jf_jb_j$ 构成的子空间 $L^p\otimes B$ 在 $1\le p<\infty$ 时稠密, 并可先在这个子空间定义 Bochner 积分, 再连续延拓到 $L^1(B)$.

若 $F\in L^p(B)$, $G\in L^{p'}(B^*)$, 则 $\langle F,G\rangle$ 可积. 总有 $L^{p'}(B^*)\subset(L^p(B))^*$; 一般不一定相等, 但当 $1<p<\infty$ 且 $B$ 自反时相等.

设 $A,B$ 是 Banach 空间, $\mathcal L(A,B)$ 是从 $A$ 到 $B$ 的有界线性算子空间. 令 $K$ 是 $\mathbb R^n\times\mathbb R^n\setminus\Delta$ 上取值于 $\mathcal L(A,B)$ 的函数, 且 $T$ 与核 $K$ 相联系, 即对具有紧支集的 $f\in L^\infty(A)$,
\[
 Tf(x)=\int_{\mathbb R^n}K(x,y)f(y)\,dy,
 \qquad x\notin\operatorname{supp}f.
\]

\hypertarget{chapter-05-page-118}{}
\begin{Theorem}
\label{thm:ch5-vector-valued-calderon-zygmund}
设对某个 $1<r<\infty$, $T:L^r(A)\to L^r(B)$ 有界, 且具有相联系的核 $K$. 若
\begin{equation}
\label{eq:ch5-vector-hormander-y}
 \int_{|x-y|>2|y-z|}\|K(x,y)-K(x,z)\|_{\mathcal L(A,B)}\,dx\le C,
\end{equation}
\begin{equation}
\label{eq:ch5-vector-hormander-x}
 \int_{|x-y|>2|x-w|}\|K(x,y)-K(w,y)\|_{\mathcal L(A,B)}\,dy\le C,
\end{equation}
则 $T:L^p(A)\to L^p(B)$ 对每个 $1<p<\infty$ 有界, 并且是弱型 $(1,1)$ 的:
\[
 |\{x\in\mathbb R^n:\|Tf(x)\|_B>\lambda\}|
 \le\frac C\lambda\|f\|_{L^1(A)}.
\]
\end{Theorem}

证明先沿用标量情形: \eqref{eq:ch5-vector-hormander-y} 与 $L^r(A)\to L^r(B)$ 有界性给出弱型 $(1,1)$, Marcinkiewicz 插值给出 $1<p<r$ 的有界性. 对 $p>r$ 需要转到伴随算子. 当 $A$ 自反时, 伴随核为 $K^*(y,x)\in\mathcal L(B^*,A^*)$, 而 \eqref{eq:ch5-vector-hormander-x} 恰给出它所需的第一个 Hörmander 条件, 所以由对偶性得到结论.

一般情形下, 先把 $T$ 限制到 $A$ 的任意有限维子空间 $A_0$. 相应核的范数不超过原核, 因而所有估计与 $A_0$ 无关. 对限制算子使用有限维对偶性, 再利用 $L^p\otimes A$ 在 $L^p(A)$ 中稠密, 即得整个空间上的结论.

\hypertarget{chapter-05-page-119}{}
证明的其余细节留给读者依照标量模型补全.
